\documentclass[11pt,a4paper,oneside,openright]{book}
\usepackage[spanish,activeacute,es-noshorthands]{babel}
\usepackage[utf8]{inputenc}
\usepackage[T1]{fontenc}
\usepackage{lmodern}
\usepackage{amsmath, amssymb, amsthm}
\usepackage{graphicx}
\usepackage[dvipsnames]{xcolor}
\usepackage{geometry}
\geometry{margin=2.5cm, headheight=14pt}
\usepackage{microtype}
\usepackage{booktabs}
\usepackage{enumitem}
\usepackage{fancyhdr}
\usepackage{tcolorbox}
\tcbuselibrary{skins, breakable}
\usepackage{caption}
\captionsetup{font=small, labelfont=bf}
\usepackage{hyperref}
\hypersetup{colorlinks=true, linkcolor=NavyBlue, urlcolor=BrickRed,
  pdftitle={M2 --- Sistemas de potencia para CubeSat}, pdfauthor={INTISAT}}

\newtheorem{definicion}{Definicion}[chapter]
\newtheorem*{nota}{Nota}

\newtcolorbox{intuibox}[1][]{colback=green!5, colframe=green!40!black,
  fonttitle=\bfseries, title=Intuicion, breakable, sharp corners=south, #1}
\newtcolorbox{calculo}[1][]{colback=gray!4, colframe=gray!55,
  fonttitle=\bfseries, title=Calculo, breakable, sharp corners=south, #1}
\newtcolorbox{payload}[1][]{colback=blue!4, colframe=blue!50!black,
  fonttitle=\bfseries, title=Conexion con el INTISAT, breakable, sharp corners=south, #1}
\newtcolorbox{cuidado}[1][]{colback=red!5, colframe=red!50!black,
  fonttitle=\bfseries, title=Cuidado, breakable, sharp corners=south, #1}
\newtcolorbox{ejercicio}[1][]{colback=orange!5, colframe=orange!60!black,
  fonttitle=\bfseries, title=Ejercicios, breakable, sharp corners=south, #1}

\pagestyle{fancy}
\fancyhf{}
\fancyhead[LE,RO]{\thepage}
\fancyhead[RE]{\nouppercase{\leftmark}}
\fancyhead[LO]{\nouppercase{\rightmark}}

\title{\Huge\bfseries M2 --- Sistemas de potencia\\[0.4em]
\Large EPS (Electrical Power System) para CubeSat\\[0.6em]
\large Manual de la coleccion CubeSat (M2 / 11)}
\author{Coleccion CubeSat para el proyecto INTISAT}
\date{\today}

\begin{document}

\frontmatter
\maketitle

\chapter*{Prefacio}
\addcontentsline{toc}{chapter}{Prefacio}

El sistema de potencia es la columna vertebral de un CubeSat. Si fallan
las comms, perdes una mision; si falla el EPS, perdes el satelite
completo. Este manual cubre como generar, almacenar y distribuir la
energia que el resto del satelite necesita.

\section*{Audiencia}
Electronica de potencia, electronica analogica, integradores de sistema.

\section*{Prerequisitos}
M0 (CubeSat 101), Manual de electronica practica (capitulos 1-8).

\tableofcontents

\mainmatter

\chapter{Generacion solar}

\section{Por que el sol es la unica opcion}

En LEO el satelite recibe energia solar 60--65\% del tiempo (el resto es
eclipse). Las opciones de energia primaria son:

\begin{itemize}[leftmargin=2em]
  \item \textbf{Solar}: lo usado en todo CubeSat. Tipica eficiencia 28\%
        (celdas triple-juncion GaAs).
  \item \textbf{Nuclear (RTG)}: caro, prohibido para mas misiones
        academicas, usado solo en deep space.
  \item \textbf{Chemical}: descartado (no se recarga).
\end{itemize}

\section{Tipos de celdas}

\begin{center}
\begin{tabular}{lll}
\toprule
\textbf{Tipo} & \textbf{Eficiencia} & \textbf{Costo relativo} \\
\midrule
Silicio (Si) cristalino       & 15-18\%  & 1x \\
Silicio monocristalino HE     & 22\%     & 1.5x \\
GaAs single-junction           & 22-24\%  & 5x \\
\textbf{Triple-junction GaAs}  & \textbf{28-32\%} & \textbf{15x} \\
Triple-junction with concentr & 35-40\%  & 30x \\
\bottomrule
\end{tabular}
\end{center}

\textbf{Para CubeSat}: casi todos usan \textbf{triple-junction GaAs}
pese al costo. La masa es limitada y la eficiencia importa mucho.

\section{Calculo de generacion}

\begin{calculo}
La constante solar a 1 AU (en LEO):
\[
S = 1361\;\text{W/m}^2
\]
Para una celda de area $A = 30\;\text{cm}^2 = 0.003\;\text{m}^2$ con
eficiencia $\eta = 0.28$:
\[
P = S \cdot A \cdot \eta = 1361 \cdot 0.003 \cdot 0.28 = 1.14\;\text{W}
\]
\end{calculo}

Para un \textbf{3U} con 4 caras laterales cubiertas (cada cara
$10 \times 30$ cm $= 300$ cm$^2$):
\begin{itemize}[leftmargin=2em]
  \item Por cara: $\sim 11$ W si recibe directo el sol.
  \item En la practica solo 1-2 caras estan iluminadas a la vez (depende
        de la orientacion).
  \item Promedio orbital con sun-pointing: $\sim 8-10$ W.
\end{itemize}

\section{Factores de degradacion}

La eficiencia BOL (Beginning of Life) cae con el tiempo:

\begin{itemize}[leftmargin=2em]
  \item \textbf{Radiacion UV}: 1-2\% por año.
  \item \textbf{Particulas energeticas}: 3-5\% por año en LEO.
  \item \textbf{Micrometeoritos}: rasguñan la superficie.
  \item \textbf{Termico}: ciclos termicos degradan los contactos.
\end{itemize}

\textbf{EOL (End of Life) tipico para 3 años en LEO}: 80-85\% del BOL.

\section{Angulo de incidencia}

La potencia generada cae con $\cos\theta$ donde $\theta$ es el angulo
entre el sol y la normal del panel:
\[
P_{\text{real}} = P_{\text{max}} \cdot \cos\theta
\]

Para mantener $\theta$ bajo, el ADCS apunta el panel al sol (sun-pointing
mode). Si el ADCS no puede, el sistema tumblea y la generacion oscila.

\begin{payload}
Para el INTISAT, asumir generacion BOL de $\sim 6-8$ W promedio orbital
(considerando que el sat tiene paneles fijos sin apuntamiento perfecto).
\end{payload}

\chapter{MPPT --- Maximum Power Point Tracking}

\section{El problema}

Una celda solar tiene una curva $V$-$I$ no lineal. La maxima potencia se
extrae a un punto especifico (MPP) que cambia con la temperatura y la
iluminacion.

Si la conectas directo a la bateria (que tiene voltaje fijo), no estas
en el MPP $\to$ perdes 20-30\% de la energia disponible.

\section{La solucion: MPPT}

Un MPPT es un \textbf{convertidor DC-DC} (buck o boost) que ajusta su
ratio para que la celda opere en el MPP. Algoritmos comunes:

\begin{itemize}[leftmargin=2em]
  \item \textbf{Perturb \& Observe (P\&O)}: cambia el ratio un poco, mide
        $P$, decide si seguir o volver.
  \item \textbf{Incremental Conductance}: mas eficiente que P\&O en
        cambios bruscos de iluminacion.
\end{itemize}

Eficiencia tipica del MPPT: 92-95\%.

\chapter{Bateria de Litio}

\section{Quimicas usadas}

\begin{center}
\begin{tabular}{lll}
\toprule
\textbf{Quimica} & \textbf{Voltaje nominal} & \textbf{Energia (Wh/kg)} \\
\midrule
Li-ion (LCO)        & 3.7 V & 150-200 \\
Li-Polymer (LiPo)   & 3.7 V & 150-200 \\
LiFePO4             & 3.2 V & 90-120 \\
NMC                 & 3.7 V & 200-220 \\
\bottomrule
\end{tabular}
\end{center}

\textbf{Para CubeSat}: la mayoria usa Li-ion 18650 estandar o Li-Polymer
custom.

\section{Voltajes clave}

\begin{itemize}[leftmargin=2em]
  \item \textbf{Cargada full}: 4.20 V (NO superar; daño permanente).
  \item \textbf{Nominal}: 3.70 V.
  \item \textbf{Descargada}: 3.0 V (NO bajar; daño irreversible).
\end{itemize}

\section{Capacidad}

Se mide en mAh o Wh:
\[
\text{Wh} = \text{V}_{\text{nominal}} \times \text{Ah}
\]

Ejemplo: bateria 18650 de 3500 mAh @ 3.7 V $=$ $3.5 \times 3.7 = 12.95$ Wh.

\begin{calculo}
Tu INTISAT 3U necesita $\sim 1.8$ Wh por orbita (estimacion).
Con bateria de $25$ Wh tenes capacidad para $25/1.8 = 14$ orbitas
sin generacion solar (solo eclipses), con DoD 100\%.
\end{calculo}

\section{Depth of Discharge (DoD)}

La bateria Li-ion sufre cuanto mas profundo la descargas:

\begin{center}
\begin{tabular}{ll}
\toprule
\textbf{DoD} & \textbf{Ciclos hasta 80\% capacidad} \\
\midrule
100\% & 500 ciclos \\
50\%  & 1500 ciclos \\
20\%  & 5000 ciclos \\
10\%  & 10000 ciclos \\
\bottomrule
\end{tabular}
\end{center}

\textbf{Para CubeSat de 1 año}: ~5500 orbitas. Si DoD nominal es 20\%,
la bateria llega al EOL con $> 90\%$ de capacidad.

\section{BMS (Battery Management System)}

Funciones del BMS:
\begin{itemize}[leftmargin=2em]
  \item Proteccion contra sobre-voltaje (cut-off > 4.2 V).
  \item Proteccion contra sub-voltaje (cut-off < 3.0 V).
  \item Proteccion contra sobre-corriente (corto-circuito).
  \item Balance entre celdas (cuando hay multiples).
  \item Monitor de temperatura.
  \item Telemetria del estado (SoC, SoH).
\end{itemize}

\chapter{Reguladores}

\section{Tipos}

\textbf{LDO (lineal)}:
\begin{itemize}[leftmargin=2em]
  \item Eficiencia: $V_{out}/V_{in}$.
  \item Disipa $(V_{in} - V_{out}) \cdot I$ como calor.
  \item Ruido bajo, simple.
  \item Mejor para diferencias pequeñas ($< 0.5$ V).
\end{itemize}

\textbf{Switching (buck/boost)}:
\begin{itemize}[leftmargin=2em]
  \item Eficiencia: $> 90\%$.
  \item Ruidoso (switching freq).
  \item Mas componentes (inductor, diodo).
  \item Mejor para conversiones grandes.
\end{itemize}

\section{Topologias usadas}

\begin{center}
\begin{tabular}{lll}
\toprule
\textbf{Topologia} & \textbf{Direccion} & \textbf{Cuando} \\
\midrule
Buck                & baja voltaje      & 5V $\to$ 3.3V \\
Boost               & sube voltaje      & 3.7V (bat) $\to$ 5V \\
Buck-boost          & sube y baja        & cuando $V_{in}$ varia \\
SEPIC                & aislado          & seguridad \\
Flyback              & aislado            & potencia mayor \\
\bottomrule
\end{tabular}
\end{center}

\chapter{Power budget}

\section{Como armar la tabla}

\begin{enumerate}[leftmargin=2em]
  \item Listar todos los subsistemas y modos.
  \item Para cada uno: potencia y duracion en orbita.
  \item Sumar $E = P \cdot t$.
  \item Comparar con generacion en orbita.
  \item Verificar margen $> 20\%$.
\end{enumerate}

\section{Ejemplo INTISAT}

\begin{center}
\small
\begin{tabular}{lcccc}
\toprule
\textbf{Modo} & \textbf{P (W)} & \textbf{t (min)} & \textbf{E (Wh)} & \textbf{\% orbita} \\
\midrule
IDLE                  & 3.0 & 70  & 3.50 & 73\% \\
CAPTURING (4 scans)   & 6.0 & 4   & 0.40 & 4\% \\
PROCESSING            & 7.5 & 12  & 1.50 & 13\% \\
DOWNLINK              & 4.0 & 8   & 0.53 & 8\% \\
SAFE (eclipse)         & 1.0 & 2   & 0.03 & 2\% \\
\midrule
\textbf{Total}        & --- & \textbf{96} & \textbf{5.96 Wh} & 100\% \\
\bottomrule
\end{tabular}
\end{center}

\textbf{Generacion estimada}: 6 W $\times$ 60 min iluminados $= 6$ Wh.
Margen: $(6 - 5.96)/6 = 0.7\%$ (muy ajustado).

\begin{cuidado}
El margen propuesto es \textbf{insuficiente}. Margen estandar para CubeSat
academico es $\geq 20\%$. Decisiones para mejorar:
\begin{itemize}[leftmargin=1em]
  \item Reducir scans a 2 por orbita.
  \item Optimizar PROCESSING (modo single en lugar de FPM).
  \item Mejorar generacion (paneles deployables).
\end{itemize}
\end{cuidado}

\chapter{EPS comerciales}

\section{Opciones de mercado}

\begin{center}
\begin{tabular}{p{4cm}p{4cm}p{2cm}}
\toprule
\textbf{Producto} & \textbf{Fabricante} & \textbf{Precio} \\
\midrule
NanoPower P31u           & GomSpace      & USD 4500 \\
iEPS-A                   & ISIS          & USD 5000 \\
EPS-1                    & EnduroSat     & USD 3500 \\
EPS Optimal              & AAC Clyde     & USD 6000 \\
Open Source EPS          & Open Cosmos   & USD 2000 \\
\bottomrule
\end{tabular}
\end{center}

Cualquier EPS comercial:
\begin{itemize}[leftmargin=2em]
  \item Acepta paneles solares.
  \item Tiene MPPT.
  \item Maneja la bateria con BMS.
  \item Genera rails 3.3V, 5V, 12V.
  \item Provee telemetria por I2C.
\end{itemize}

\textbf{Para INTISAT}: idealmente un EPS comercial probado, no diseñarlo
desde cero (alto riesgo).

\chapter{Test de potencia}

\section{Tests en banco}

\begin{itemize}[leftmargin=2em]
  \item \textbf{Caracterizacion de panel solar}: bajo lampara halogena
        simulada, medir curva I-V.
  \item \textbf{Test del MPPT}: variar irradiancia, verificar que
        rastrea el MPP.
  \item \textbf{Test de BMS}: descargar bateria hasta cutoff, verificar
        proteccion.
  \item \textbf{Test de regulacion}: variar carga, medir ripple.
  \item \textbf{Test de eficiencia}: medir eficiencia global del sistema.
\end{itemize}

\section{Tests ambientales}

Para el EPS especificamente:
\begin{itemize}[leftmargin=2em]
  \item \textbf{TVAC}: ciclos -20 a +60 °C. Verificar que la bateria
        funciona en frio.
  \item \textbf{Vibracion}: el packaging de las baterias 18650 puede
        soltarse.
  \item \textbf{Vacio}: bateria Li-ion no deberia inflarse en vacio
        (electrolito gaseoso).
\end{itemize}

\backmatter

\chapter*{Bibliografia}
\addcontentsline{toc}{chapter}{Bibliografia}

\begin{enumerate}[leftmargin=2em]
\item Larson WJ, Wertz JR. \emph{Space Mission Analysis and Design}. 3ed. 1999.
\item Patel MR. \emph{Spacecraft Power Systems}. CRC Press, 2005.
\item GomSpace. \emph{NanoPower P31u Datasheet}. \url{https://gomspace.com}.
\item ISIS. \emph{iEPS Datasheet}. \url{https://www.isispace.nl}.
\item Bouwmeester J. State of the art of CubeSats. \emph{Acta Astronautica}, 2010.
\end{enumerate}

\end{document}
